%%%%%%%%%%%%%%%%%%%%%%%%%%%%%%%%%%%%%%%%%%%%%%%%%%%%%%%%%%%%%%%%%%%%%%%%
\chapter{Applicazioni della CP}
%%%%%%%%%%%%%%%%%%%%%%%%%%%%%%%%%%%%%%%%%%%%%%%%%%%%%%%%%%%%%%%%%%%%%%%%

\section{Rappresentare forme bilineari tramite tensori}
\label{sec:tensor_bil}

Conoscere il rango CP di tensori che rappresentano applicazioni multilineari può rivelare
informazioni sulla complessità asintotica e fornire algoritmi con complessità ottimale per la loro
valutazione. Proviamo a formulare il problema: dall'algebra lineare sappiamo che data una forma
bilineare $ b : \R^m \times \R^n \to \R $, esiste una matrice $ M \in \R^{m \times n} $ tale che $
b(u,v) = u^{\top}Mv$ per ogni $ (u,v) \in \R^m \times \R^n $.
Notiamo che un insieme di $ p $ forme bilineari, o equivalentemente una funzione \todo{È chiaro se
si chiama ancora b?}bilineare $
b = \begin{bmatrix} b_1 \\ \vdots \\ b_p \end{bmatrix} :  \R^m \times \R^n  \to \R^p$ si può
rappresentare tramite un tensore $ \Tn{x} \in \R^{m \times n\times p} $ tale che le sue slice
frontali rappresentino ognuna delle $ p $ forme bilineari, ossia che
\begin{equation}
  \label{eq:b_k_bil}
  b_k(u,v) = u^{\top}\Tn{x}(:,:,k)v = \sum_{i=1}^m\sum_{j=1}^n \Tn{x}(i,j,k)u_iv_j, \text{ per ogni
  } k = 1, \ldots, p
.\end{equation}
Da cui, 
\begin{equation}
  \label{eq:b_tuv}
 b(u,v) = \begin{bmatrix} b_1(u,v) \\ \vdots \\ b_p(u,v) \end{bmatrix} = \Tn{x} \ttm[1] u^{\top}
 \ttm[2] v^{\top}.
\end{equation}
\begin{remark}
  Questa procedura si estende anche al caso in cui le entrate di $ u $ e $ v $ siano oggetti
  algebrici più complessi di scalari, ad esempio vettori o matrici.
\end{remark}
\begin{remark}[Active multiplications]
  Scrivere una forma bilineare come nell'Equazione \ref{eq:b_k_bil} evidenzia il numero di
  moltiplicazioni necessarie fra gli elementi di $ u $ e $ v $. Queste sono dette \emph{active
  multiplications} e corrispondono al numero di elementi non nulli di $ \Tn{x} $. Poiché ci aspettiamo che
  il costo delle active multiplications sia quello che determini la complessità di valutare $ b(u,v)
  $, cerchiamo di minimizzarne il numero cercando fattorizzazioni CP di $ \Tn{x} $.
\end{remark}
Se abbiamo $ \Tn{x} = \dsqr{A,B,C} $ con $ \rk(\Tn{x}) = r$, dalla Definizione~\ref{def:CP} abbiamo che $ \Tn{x} = \sum_{l=1}^r a_l \topp b_l \topp c_l $, o in componenti, $ \Tn{x}(i,j,k) = \sum_{l=1}^r a_{il}b_{jl}c_{kl} $. Inserendo questa relazione nella Equazione~\ref{eq:b_k_bil} troviamo,

\begin{equation}
  \label{eq:b_active}
  \begin{aligned}
    b_k(u,v) &= \left( \left( \sum_{i=1}^r a_l \topp b_l \topp c_l \right) \ttm[1] u^{\top} \ttm[2]
  v^{\top}  \right)_k  
             = \left( \sum_{l=1}^r \left( a_l \topp b_l \topp c_l  \right) \ttm[1] u^{\top} \ttm[2]
    v^{\top} \right)_k  \\ 
             &=
  \sum_{i=1}^m \sum_{j=1}^n \sum_{l=1}^r a_{il}b_{jl}c_{kl} u_i v_j = \sum_{l=1}^r \left(\sum_{i=1}^m 
   a_{il} u_i \right) \cdot \left( \sum_{j=1}^n b_{jl} v_j  \right)  c_{kl} = \sum_{l=1}^r
   (a_l^{\top}u) \cdot (b_l^{\top}v)c_{kl},
  \end{aligned}
\end{equation}
per ogni $ k = 1, \ldots, p $.
Mettendo insieme le Equazioni \ref{eq:b_tuv} e \ref{eq:b_active} troviamo
\begin{equation*}
  b(u,v) = \sum_{l=1}^r (a_l^{\top}u)\cdot(b_l^{\top}v)c_l.
\end{equation*}

Quindi valutare $ b $ richiede esattamente $ r $ active multiplications, che evidenziamo con $
(\cdot) $.
\todo{Espandere discussione dall'articolo}


\section{Il master method}



Introduciamo brevemente un metodo per risolvere relazioni ricorsive definite per ricorrenza, detto
\emph{master method}. I lettori più interessati possono trovare una trattazione più dettagliata sul
libro~\cite{cormen2022introduction}. Il master method si usa per risolvere ricorrenze della forma
\begin{equation}
  \label{eq:master_rec}
  T(n) = aT(\frac{n}{b}) + f(n),
\end{equation}
dove $ a \ge 1 $ e $ b>1 $ sono costanti e $ f(n) $ è una funzione asintoticamente positiva. La
ricorrenza \ref{eq:master_rec} descrive il tempo di esecuzione di un algoritmo che divide un
problema di dimensione $ n $ in $ a $ sottoproblemi, ognuno di dimensione $ \frac{n}{b} $. Gli $ a $
sottoproblemi vengono risolti ricorsivamente in tempo $ T(\frac{n}{b}) $, e la funzione $ f(n) $
rappresenta il costo di dividere il problema e ricombinare la ricorsione di ogni sottoproblema.

Indichiamo con $ O $ e $ \Omega $ la stima superiore e inferiore, rispettivamente, e con $
\Theta $ la stima stretta.
\begin{theorem}[Master theorem]
  \label{th:master}
  Siano $ a\ge 1 $ e $ b>1 $ delle costanti, sia $ f(n) $ una funzione e definiamo $ T(n) $ sugli
  interi nonnegativi dalla ricorrenza
  \begin{equation*}
    T(n) = aT(\frac{n}{b}) + f(n),
  \end{equation*}
\end{theorem}
dove con $ \frac{n}{b} $ indichiamo sia $ \left\lfloor \frac{n}{b} \right\rfloor $ o $ \left\lceil
\frac{n}{b} \right\rceil  $. Allora $ T(n) $ ha le seguenti soglie asintotiche:
\begin{enumerate}
  \item Se $ f(n) = O(n^{\log_b a-\varepsilon}) $ per qualche $ \varepsilon >0 $, allora $ T(n) =
    \Theta(n^{\log_b a}) $.
  \item Se $ f(n) = \Theta(n^{\log_b a}) $ allora $ T(n) = \Theta(n^{\log_b a \log n}) $.
  \item If $ f(n) = \Omega(n^{\log_b a + \varepsilon}) $ per qualche costante $ \varepsilon >0 $, e
    se $ af(\frac{n}{b}) < cf(n) $ per qualche costante $ c<1$ e ogni $ n $ sufficientemente grande,
    allora $ T(n) = \Theta(f(n)) $.
\end{enumerate}

\begin{example}
  Nel caso dell'algoritmo naive la ricorrenza ha la forma $ a = 8 $, $ b = 2 $ e $ f(n) =
  \Theta(n^2) $, ovvero
  \begin{equation*}
    T(n) = 8T(\frac{n}{2}) + \Theta(n^2).
  \end{equation*}
  Siccome $ n^{\log_b a} = n^{\log_2 8} = n^3 $, che è polinomialmente più grande di $ f(n) $, ovvero $
  f(n) = O(n^{3-\varepsilon}) $ per $ \varepsilon = 1 $, dal caso 1 di \ref{th:master} segue che $
  T(n) = \Theta(n^3) $.
\end{example}


\section{Moltiplicazione ricorsiva di matrici}

Possiamo rappresentare la moltiplicazione tra matrici come una funzione bilineare, quindi come un
tensore. 

\begin{lemma}[Rappresentazione del prodotto matriciale]
Per ogni coppia di matrici $ A \in \R^{m \times n} $ e $ B \in \R^{n
\times p} $, esiste un tensore $ \Tn{x} \in ? $ tale che il prodotto $ C = AB \in \R^{m \times n} $
si scrive come $ \Tn{x} \ttm[1] \vec(A) \ttm[2] \vec(B) =
\vec(C)$.
\end{lemma}
\todo{finire}
\begin{proof}
  Dalla definizione di prodotto di matrice, $ C_{ij} =  $
\end{proof}

\todo{Brief motivation}

\begin{notation}
  Indichiamo con $ \angb{m,n,p} $ il prodotto tra due matrici di taglia $ m \times n $ e $ n \times
  p $.
\end{notation}


Trattiamo prima il caso di matrici quadrate con taglie potenze di due. Discuteremo di seguito come
trattare il caso generale. In questo caso possiamo
partizionare le matrici in delle sottomatrici a blocchi $ 2\times 2 $:
\begin{equation*}
  \begin{bmatrix}
    C_{11} & C_{12} \\
    C_{21} & C_{22}
  \end{bmatrix} =
  \begin{bmatrix}
    A_{11} & A_{12} \\
    A_{21} & A_{22}
  \end{bmatrix} \cdot 
  \begin{bmatrix}
    B_{11} & B_{12} \\
    B_{21} & B_{22}
  \end{bmatrix} =
  \begin{bmatrix}
    A_{11}B_{11} + A_{12}B_{21} & A_{11}B_{12} + A_{12}B_{22} \\
    A_{21}B_{12} + A_{22}B_{21} & A_{21}B_{12} + A_{22}B_{22}
  \end{bmatrix}.
\end{equation*}
L'algoritmo naive procede combinando un insieme di otto moltiplicazioni matriciali con quattro
addizioni:
% Intermediate Matrices
% Intermediate Matrices
\[
\begin{aligned}
M_1 &= A_{11} \cdot B_{11} &\quad M_2 &= A_{12} \cdot B_{21} &\quad M_3 &= A_{11} \cdot B_{12} \\
M_4 &= A_{12} \cdot B_{22} &\quad M_5 &= A_{21} \cdot B_{11} &\quad M_6 &= A_{22} \cdot B_{21} \\
M_7 &= A_{21} \cdot B_{12} &\quad M_8 &= A_{22} \cdot B_{22} & &
\end{aligned}
\]

% Resulting Matrices
\[
\begin{aligned}
C_{11} &= M_1 + M_2 &\qquad C_{12} &= M_3 + M_4 \\
C_{21} &= M_5 + M_6 &\qquad C_{22} &= M_7 + M_8
\end{aligned}
\]
